% !TEX recipe = pdflatex
% Nova apresentacao; modelo visual: Apresentacao 24092026 - Hipoteses e RL/topicos_micro.tex
\documentclass[aspectratio=169,11pt]{beamer}
\usetheme{Madrid}
\usepackage[T1]{fontenc}
\usepackage[utf8]{inputenc}
\usepackage[brazilian]{babel}
\usepackage{lmodern}
\usepackage{booktabs,array,tabularx}
\definecolor{verde}{HTML}{24574D}
\definecolor{grafite}{HTML}{253A40}
\setbeamercolor{structure}{fg=verde}
\setbeamercolor{title}{bg=verde,fg=white}
\setbeamercolor{frametitle}{bg=verde,fg=white}
\setbeamercolor{normal text}{fg=grafite}
\setbeamertemplate{navigation symbols}{}
\setbeamertemplate{footline}{\leavevmode\hbox{\begin{beamercolorbox}[wd=.86\paperwidth,ht=2.7ex,dp=1.2ex,leftskip=1em]{author in head/foot}Capital natural e governança territorial\end{beamercolorbox}\begin{beamercolorbox}[wd=.14\paperwidth,ht=2.7ex,dp=1.2ex,center]{date in head/foot}\insertframenumber\,/\,\inserttotalframenumber\end{beamercolorbox}}}
\setbeamersize{text margin left=8mm,text margin right=8mm}
\setbeamertemplate{itemize items}[circle]
\newcommand{\fonte}[1]{\par\vfill{\fontsize{7.5}{9}\selectfont\color{gray}#1\par}}
\newcommand{\tpar}{\par\vspace{1.5mm}}
\newcolumntype{P}[1]{>{\raggedright\arraybackslash}p{#1}}
\newcommand{\disp}{\textcolor{verde}{Disponível}}
\newcommand{\parcial}{\textcolor{orange!75!black}{Parcial}}
\newcommand{\pend}{\textcolor{red!70!black}{Pendente}}
\title[Capital natural e governança territorial]{Capital Natural e Governança Territorial}
\subtitle{Hipóteses, dados e estratégia econométrica}
\author{Fernando Souza da Vieira}
\institute{Tópicos de Microeconomia}
\date{}
\begin{document}

\begin{frame}
  \titlepage
\end{frame}

\begin{frame}{Problema de pesquisa em Mato Grosso}
\begin{block}{Pergunta central}
Como retornos econômicos e instituições condicionam as decisões de conversão ou manutenção da vegetação, e em que medida o estoque de vegetação nativa altera a resposta da produtividade agrícola a choques climáticos em contextos ecologicamente e economicamente heterogêneos?
\end{block}
\medskip
\begin{columns}[T]
\column{.49\textwidth}
\textbf{Decisão de uso da terra}\par\smallskip
Retornos, custos, riscos e instituições se relacionam com conversão ou conservação.
\column{.49\textwidth}
\textbf{Resposta produtiva}\par\smallskip
O estoque vegetal prévio pode alterar o efeito do choque climático sobre a produtividade, com diferenças entre biomas.
\end{columns}
\medskip
\centering\small\color{verde}
Retornos + instituições $\rightarrow$ conversão $\rightarrow$ vegetação\par
Vegetação prévia $\times$ choque climático $\rightarrow$ produtividade
\fonte{Pergunta e duas dimensões: introdução atual, linhas 81--87. A retroalimentação permanece mecanismo conceitual.}
\end{frame}

\begin{frame}{H1: retornos, preços e instituições}
\begin{block}{H1 — Retornos, instituições e conversão}
A taxa de conversão de vegetação nativa para usos agropecuários aumenta quando se eleva o retorno econômico esperado desses usos.
\end{block}
\smallskip
\begin{block}{H1-a — Preços e capacidade produtiva local}
O efeito de aumentos nos preços dos produtos agropecuários sobre a taxa de conversão de vegetação nativa é maior onde a capacidade produtiva local da terra para esses usos é mais elevada.
\end{block}
\smallskip
\begin{block}{H1-b — Condicionamento institucional}
A resposta da taxa de conversão ao retorno econômico esperado dos usos agropecuários é condicionada pelo ambiente institucional, particularmente pelas regras, pelo monitoramento, pela fiscalização e pelas sanções que alteram os custos e os riscos associados à conversão.
\end{block}
\end{frame}

\begin{frame}{H1: desfecho e capacidade produtiva}
\fontsize{8.8}{10.3}\selectfont
\setlength{\tabcolsep}{3pt}
\begin{tabularx}{\textwidth}{P{2.0cm} P{3.45cm} P{2.0cm} P{1.8cm} X}
\toprule
\textbf{Variável / insumo} & \textbf{Descrição da variável} & \textbf{Fonte} & \textbf{Período} & \textbf{Situação} \\
\midrule
$Conv_{it}$, $Conv^{Agric}_{it}$, $Conv^{Past}_{it}$ & Área nativa em $t-1$ convertida em $t$ / vegetação nativa disponível em $t-1$; destinos agropecuários. & MapBiomas 10.1 & 1985--2024 & \pend: transições brutas; variação líquida não substitui conversão. \\
Produção agrícola & Área, produção e rendimento de soja e milho; apoio à caracterização produtiva. & IBGE/PAM & 1974--2025 & \disp; cobertura municipal variável. \\
Pecuária & Efetivo bovino; indicador da atividade pecuária, não da conversão em pastagem. & IBGE/PPM & 1974--2024 & \disp; 2025 não retornou. \\
$S_{ik}$ & Aptidão ou rendimento atingível candidato para soja e milho; não é rendimento observado. & FAO/GAEZ v4 & clima 1961--1990 & \parcial; base de terras inclui 2010. \\
\bottomrule
\end{tabularx}
\end{frame}

\begin{frame}{H1: preços e instituições}
\fontsize{8.8}{10.3}\selectfont
\setlength{\tabcolsep}{3pt}
\begin{tabularx}{\textwidth}{P{2.0cm} P{3.45cm} P{2.0cm} P{1.8cm} X}
\toprule
\textbf{Variável / insumo} & \textbf{Descrição da variável} & \textbf{Fonte} & \textbf{Período} & \textbf{Situação} \\
\midrule
$P_{kt}$ & Preços mensais de soja, milho e carne bovina; determinantes candidatos do retorno. & World Bank Pink Sheet & jan. 1960--ago. 2026 & \disp; USD nominais. \\
Preços e custos locais & Preços e custos para aproximações alternativas do retorno. & CEPEA / IMEA / CONAB & séries não consolidadas & \pend; custos CONAB extraídos. \\
$I_{i,t-1}$ & Autos e multas candidatos; não equivalem a sanções efetivas. & IBAMA & 1987--2026 & \parcial; 52.991 eventos, 2026 parcial. \\
Regras / cadastro & Vigência normativa e informação fundiária para caracterizar instituições. & MMA /\newline CAR & histórico insuficiente & \pend. \\
\bottomrule
\end{tabularx}
\end{frame}

\begin{frame}{H2: vegetação, clima e produtividade}
\begin{block}{H2 — Vegetação, clima e produtividade}
A resposta da produtividade agrícola aos choques climáticos varia, em magnitude economicamente relevante, conforme o estoque de vegetação previamente existente no território.
\end{block}
\smallskip
\begin{block}{H2-a — Heterogeneidade entre biomas}
A magnitude da modificação exercida pelo estoque de vegetação sobre a resposta da produtividade agrícola aos choques climáticos é heterogênea entre biomas.
\end{block}
\medskip
\begin{block}{Variáveis candidatas}
$Y_{it}$: rendimento da soja; $C_{it}$: choque climático; $E_{i,t-1}$: vegetação prévia. O interesse recai sobre $C_{it}\times E_{i,t-1}$ e sua heterogeneidade entre biomas.
\end{block}
\fonte{Texto integral das hipóteses da introdução atual. H2 não impõe sinal nem maior produtividade média.}
\end{frame}

\begin{frame}{H2: dados, fontes e períodos}
\fontsize{8.5}{10.2}\selectfont
\setlength{\tabcolsep}{3.5pt}
\begin{tabularx}{\textwidth}{P{2.1cm} P{2.85cm} P{2.2cm} X}
\toprule
\textbf{Variável / insumo} & \textbf{Fonte} & \textbf{Cobertura da fonte} & \textbf{Situação para H2} \\
\midrule
Rendimento da soja & IBGE/PAM & 1974--2025 & \disp; símbolos SIDRA e cobertura por município a revisar. \\
Vegetação prévia & MapBiomas 10.1 & 1985--2024 & \parcial: estoque candidato; taxonomia e denominador pendentes. \\
Chuva / anomalias & CHIRPS v2 & 1981--2025; 540 meses & \parcial: chuva, déficit, excesso e z-score em ano civil; janela agrícola pendente. \\
Temperatura & ERA5-Land & não adquirida & \pend: configuração e condições de acesso CDS. \\
Biomas & IBGE 2019 / malha 2024 & corte espacial & \parcial: limiar de 80\% candidato; sensibilidades de 70\% e 90\%. \\
Topografia & Copernicus GLO-90 & corte espacial & \disp: elevação de modelo de superfície; declividade não calculada. \\
Calendário agrícola & ZARC/MAPA/\newline CONAB & safra 2025/26 & \parcial: documento não cobre calendário histórico. \\
\bottomrule
\end{tabularx}
\smallskip
\textbf{1986--2024: janela potencial} de vegetação defasada e chuva, não amostra final. Faltam geografia histórica, janela agrícola e choque definitivo.
\fonte{Fonte: auditoria de 30/09/2026. Temperatura depende do desenho escolhido.}
\end{frame}

\begin{frame}{Modelos econométricos candidatos}
\fontsize{9}{11}\selectfont
\setlength{\tabcolsep}{4pt}
\begin{tabularx}{\textwidth}{P{1.2cm} P{2.25cm} P{4.4cm} X}
\toprule
\textbf{Hipótese} & \textbf{Dependente} & \textbf{Independentes centrais} & \textbf{Modelo candidato} \\
\midrule
H1 & $Conv_{it}$ & Retorno esperado $R^e_{i,t-1}$ ou determinantes observáveis. & Painel municipal com efeitos fixos de município e tempo. \\
H1-a & $Conv_{it}$ & Choques de preços $\Delta\ln P_{k,t-1}$ interagidos com capacidade $S_{ik}$. & Painel com interação preço $\times$ capacidade. \\
H1-b & $Conv_{it}$ & Retorno esperado, instituição $I_{i,t-1}$ e sua interação. & Painel com interação retorno $\times$ instituição. \\
H2 & $Y_{it}$, soja kg/ha & Choque climático $C_{it}$, vegetação prévia $E_{i,t-1}$ e interação. & Painel com efeitos fixos e moderação climática. \\
H2-a & $Y_{it}$, soja kg/ha & $C_{it}\times E_{i,t-1}\times Bioma_i$. & Interação tripla ou painéis por bioma, se houver suporte. \\
\bottomrule
\end{tabularx}
\medskip
\begin{block}{Ponto de partida, sujeito à validação}
Controles, defasagens e amostra dependerão da integração dos dados. Erros agrupados por município são candidatos; desfechos proporcionais e dependência espacial exigem avaliação.
\end{block}
\fonte{Associação não implica causalidade. Fiscalização pode reagir à conversão; a identificação institucional permanece aberta.}
\end{frame}

\begin{frame}{Próximas etapas}
\begin{block}{Dados ainda não formam uma amostra econométrica}
1986--2024 é janela potencial de vegetação defasada e chuva. A interseção das variáveis, a unidade territorial e as medidas finais ainda precisam ser definidas.
\end{block}
\medskip
\begin{enumerate}
\item Harmonizar municípios e construir as conversões brutas MapBiomas.
\item Definir janela agrícola e choque climático; consolidar preços, aptidão e instituições.
\item Integrar o painel e verificar cobertura, ausências e suporte por bioma.
\item Revisar a identificação antes de definir modelos e robustez.
\end{enumerate}
\medskip
\begin{alertblock}{Escopo da apresentação}
Os modelos apresentados são candidatos; a identificação final dependerá das propriedades, limitações e variação efetivamente disponíveis nos dados.
\end{alertblock}
\end{frame}
\end{document}
